%% SYNTHETIC TEST UNIT generated by writing/tools/make_outline_book.py from FINAL_OUTLINE_v2 -- not book text
\chapter{Sudakov 形状因子与振幅的红外结构}\label{ch:sudakov}
本单元为测试用的合成文本（由 FINAL\_OUTLINE\_v2 自动生成），陈述内容从略。

\section{第 1 节（测试）}\label{sec:sudakov-oneloop}

\begin{definition}[夸克形状因子]\label{def:sudakov-ff}
（测试用合成陈述：内容从略。）
\uses{def:prelim-feynman-rules}
\end{definition}

\begin{proposition}[一圈形状因子（精确到 $\varepsilon$ 的所有阶） \tagR]\label{prop:sudakov-oneloop}
（$R_{\mathrm{PT}}$：在重整化微扰论中对红外正规化的部分子态严格。）
（测试用合成陈述：内容从略。）
（本命题只涉及部分子态。）
\uses{def:sudakov-ff, def:qcd-dimreg, def:prelim-dr-conventions, lem:prelim-scaleless, ex:prelim-scaleless}
\end{proposition}

\begin{example}[在壳无质量形状因子 = 硬区]\label{ex:sudakov-onshell-hard}
（测试用合成陈述：内容从略。）
\uses{prop:sudakov-oneloop, def:ir-regions}
\end{example}

\section{第 2 节（测试）}\label{sec:sudakov-regions}

\begin{example}[离壳 Sudakov 三角形的区域展开]\label{ex:sudakov-offshell}
（测试用合成陈述：内容从略。）
\uses{def:ir-regions, prop:ir-regions-minkowski}
\end{example}

\begin{example}[有质量交换：需要解析正规化]\label{ex:sudakov-massive}
（测试用合成陈述：内容从略。）
\uses{def:ir-regions}
\end{example}

\begin{example}[双对数系数]\label{ex:sudakov-ll}
（测试用合成陈述：内容从略。）
\uses{ex:sudakov-offshell, ex:sudakov-massive}
\end{example}

\section{第 3 节（测试）}\label{sec:sudakov-factorization}

\begin{remark}[一致性论证：单标度结构、$\mu$ 无关性与极点相消]\label{rem:sudakov-consistency}
（测试用合成陈述：内容从略。）
\uses{ex:sudakov-offshell}
\end{remark}

\begin{theorem}[Sudakov 形状因子的全阶因子化 \tagA]\label{thm:sudakov-factorization}
（测试用合成陈述：内容从略。）
（本命题只涉及部分子态。）
\uses{thm:ir-leading-regions, prop:ir-power-counting, lem:prelim-eikonal-identity, prop:prelim-wilson-rules, def:prelim-wilson, thm:prelim-slavnov-taylor, prop:sudakov-oneloop, ex:sudakov-offshell}
\end{theorem}
\begin{proofsketch}测试。\end{proofsketch}
\paragraph{地位与缺口}测试：本陈述的地位由\cref{prop:appN-triangle-cs}、\cref{asm:appN-collins}、\cref{thm:appN-subtraction}、\cref{def:method-collins-conditions}决定。

\begin{definition}[cusp 反常维度]\label{def:sudakov-cusp}
（测试用合成陈述：内容从略。）
\uses{def:prelim-wilson}
\end{definition}

\section{第 4 节（测试）}\label{sec:sudakov-evolution}

\begin{proposition}[$K+G$ 演化方程 \tagA]\label{prop:sudakov-kg}
（测试用合成陈述：内容从略。）
\uses{thm:sudakov-factorization, def:sudakov-cusp, prop:qcd-rge, prop:sudakov-oneloop}
\end{proposition}
\begin{proofsketch}测试。\end{proofsketch}
\paragraph{地位与缺口}测试：本陈述的地位由\cref{thm:sudakov-factorization}决定。

\begin{corollary}[指数化与解析延拓 \tagA]\label{cor:sudakov-exponentiation}
（测试用合成陈述：内容从略。）
\uses{prop:sudakov-kg}
\end{corollary}
\begin{proofsketch}测试。\end{proofsketch}
\paragraph{地位与缺口}测试：本陈述的地位由\cref{prop:sudakov-kg}决定。

\begin{proposition}[非阿贝尔指数化（webs） \tagA]\label{prop:sudakov-webs}
（测试用合成陈述：内容从略。）
\uses{def:prelim-wilson, prop:prelim-wilson-rules, lem:prelim-eikonal-identity, prop:appB-webs}
\end{proposition}
\begin{proofsketch}测试。\end{proofsketch}
\paragraph{地位与缺口}测试：见文献。

\section{第 5 节（测试）}\label{sec:sudakov-multileg}

\begin{theorem}[无质量多腿振幅红外极点的因子化 \tagA]\label{thm:sudakov-ir-poles}
（测试用合成陈述：内容从略。）
（本命题只涉及部分子态。）
\uses{thm:sudakov-factorization, def:sudakov-cusp, prop:sudakov-webs}
\end{theorem}
\begin{proofsketch}测试。\end{proofsketch}
\paragraph{地位与缺口}测试：本陈述的地位由\cref{prop:sudakov-dipole-violation}决定。

\begin{proposition}[偶极子公式作为全阶精确公式的失效 \tagV]\label{prop:sudakov-dipole-violation}
（测试用合成陈述：内容从略。）
\uses{thm:sudakov-ir-poles}
\end{proposition}
\paragraph{地位与缺口}测试：见文献。

\begin{proposition}[Coulomb 相位 \tagR]\label{prop:sudakov-coulomb}
（$R_{\mathrm{PT}}$：在重整化微扰论中对红外正规化的部分子态严格。）
（测试用合成陈述：内容从略。）
\uses{prop:sudakov-oneloop}
\end{proposition}

\section{第 6 节（测试）}\label{sec:sudakov-collinear}

\begin{proposition}[振幅的硬—软—共线因子化（出射或阈值运动学） \tagA]\label{prop:sudakov-amplitude-factorization}
（测试用合成陈述：内容从略。）
\uses{thm:sudakov-factorization, thm:ir-leading-regions}
\end{proposition}
\begin{proofsketch}测试。\end{proofsketch}
\paragraph{地位与缺口}测试：见文献。

\begin{remark}[类空共线极限预览 \tagV]\label{rem:sudakov-spacelike}
（测试用合成陈述：内容从略。）
\uses{prop:sudakov-amplitude-factorization}
\end{remark}
\paragraph{地位与缺口}测试：见文献。
